\documentclass[11pt,article,oneside]{memoir}
\usepackage[utf8]{inputenc}
\usepackage[T1]{fontenc}
\usepackage{url}
\usepackage[hidelinks,colorlinks=true,linkcolor=blue]{hyperref}
\usepackage{amsmath,amssymb}
\usepackage{graphicx}
\usepackage{xcolor}
\usepackage{tabularx}
\usepackage{listings}
\usepackage{geometry}
\usepackage{titlesec}
\usepackage{libertine}

\geometry{margin=1in}

\titleformat{\section}{\Large\bfseries\color{darkgray}}{\thesection}{1em}{}
\titleformat{\subsection}{\large\bfseries\color{blue!80!black}}{\thesubsection}{1em}{}

\lstset{
  basicstyle=\small\ttfamily,
  breaklines=true,
  frame=lines,
  backgroundcolor=\color{gray!10}
}

\begin{document}

\begin{center}
    \Huge \textbf{Integration Guide} \\
    \vspace{0.5em}
    \large \textit{A CEREBRUM Research Publication (v1.2.4)} \\
    \vspace{1em}
    \large \textbf{Bryan Alexander Buchorn} \\
    Independent Researcher \\
    \vspace{2em}
\end{center}

\begin{abstract}
\textbf{Overview:} This guide provides a conceptual entry point into the CEREBRUM framework. It is designed for practitioners and researchers seeking to understand the core reasoning signals of Framework Implementation Guide.
\end{abstract}

\section{CEREBRUM Integration Guide}

\textbf{Version}: v1.1.0
\textbf{Audience}: Python developers integrating CEREBRUM into an application

This guide covers the Python API directly. For the REST API, see \texttt{docs/API\_REFERENCE.md}. For deployment, see \texttt{docs/DEPLOYMENT.md}.

---

\subsection{Instal\-lation}

\begin{lstlisting}
# Core only (no sentence-transformers, no API server)
pip install -e "."

# With semantic embeddings
pip install -e ".[embeddings]"

# With REST API server
pip install -e ".[api]"

# Everything
pip install -e ".[all]"
\end{lstlisting}

---

\subsection{1. Loading a Graph}

\subsubsection{From CSV}
\begin{lstlisting}
from adapters.csv_adapter import load_csv_adapter

adapter = load_csv_adapter("path/to/graph.csv")
# CSV format: subject,predicate,object[,weight]
# Header row optional; auto-detected
\end{lstlisting}

\subsubsection{From NetworkX}
\begin{lstlisting}
import networkx as nx
from adapters.networkx_adapter import NetworkXAdapter

G = nx.DiGraph()
G.add_edge("Marie Curie", "Radium", relation="discovered", weight=0.95)
G.add_edge("Radium", "Radioactivity", relation="property_of", weight=0.90)

adapter = NetworkXAdapter(G)
\end{lstlisting}

\subsubsection{From Neo4j}
\begin{lstlisting}
from adapters.neo4j_adapter import Neo4jAdapter

adapter = Neo4jAdapter(
    uri="bolt://localhost:7687",
    user="neo4j",
    password=os.environ["NEO4J_PASSWORD"],
)
\end{lstlisting}

\subsubsection{From RDF / SPARQL}
\begin{lstlisting}
from adapters.rdf_adapter import RDFAdapter

adapter = RDFAdapter(sparql_endpoint="https://dbpedia.org/sparql")
\end{lstlisting}

\subsubsection{Progra\-mmatic const\-ruction with IngestionPipeline}
\begin{lstlisting}
from core.thalamus import IngestionPipeline

pipeline = IngestionPipeline(
    adapter,
    namespace="text",          # Phase 19: isolate from signal entities
    entity_dedup_map={         # alias normalization
        "M. Curie": "Marie Curie",
        "Madame Curie": "Marie Curie",
    }
)

pipeline.process("Marie Curie", "discovered", "Radium", confidence=0.99, source="wikipedia")
pipeline.process("Radium", "causes", "Radiation_Sickness", confidence=0.85, source="pubmed")
\end{lstlisting}

---

\subsection{2. Community Detection (DSCF)}

\begin{lstlisting}
from core.community_engine import best_of_n_dscf, modularity_score

G = adapter.to_networkx()

# Best of N trials (recommended: n_trials=5 for reproducibility)
partitions = best_of_n_dscf(G, n_trials=5, seed=42)
q = modularity_score(G, partitions)
print(f"{len(partitions)} communities, Q={q:.4f}")

# Build community_map: {node_id: community_id}
community_map = {}
for cid, members in enumerate(partitions):
    for node in members:
        community_map[node] = cid
\end{lstlisting}

\subsubsection{Algorithm selection}
\begin{lstlisting}
from core.community_engine import CommunityEngine

# DSCF (default, recommended)
engine = CommunityEngine(adapter, algorithm="dscf")

# Leiden (native, GPL-free)
engine = CommunityEngine(adapter, algorithm="leiden")

# LPA only (fast, lower quality)
engine = CommunityEngine(adapter, algorithm="lpa")

# With soft membership (Phase 17)
engine = CommunityEngine(adapter, algorithm="dscf", soft_membership=True, tau=2.0)
\end{lstlisting}

---

\subsection{3. Embeddings}

\begin{lstlisting}
from core.embedding_engine import RandomEngine, SentenceEngine

# Random (fast, no external deps, good for structural reasoning)
engine = RandomEngine(dim=64, seed=42)
embeddings = engine.encode_entities({node: node for node in G.nodes()})

# Semantic (requires: pip install -e ".[embeddings]")
engine = SentenceEngine(model="all-MiniLM-L6-v2")
embeddings = engine.encode_entities({node: node for node in G.nodes()})

# KGE embeddings (Phase 17, requires training)
from core.kge_engine import RotatEEngine, KGEEmbeddingAdapter
kge = RotatEEngine(adapter, dim=128, epochs=200)
kge.train()
embeddings = KGEEmbeddingAdapter(kge)
\end{lstlisting}

---

\subsection{4. Building the CSA Engine}

\begin{lstlisting}
from core.attention_engine import CSAEngine
from core.structural_encoder import build_community_distance_matrix, adjacent_community_pairs

dist = build_community_distance_matrix(G, community_map)
adj  = adjacent_community_pairs(G, community_map)

csa = CSAEngine(
    communities=community_map,
    embeddings=embeddings,
    # Optional: per-community parameter overrides (Phase 20)
    community_params={
        3: (0.5, 0.15, 0.25, 0.05, 0.05, 0.0),  # protein community — reduce β, increase γ
    },
)
csa.set_community_graph(dist, adj)

# Phase 20: snapshot isolation for concurrent use
csa.set_query_snapshot(community_map)
\end{lstlisting}

---

\subsection{5. Running a Query}

\begin{lstlisting}
from reasoning.traversal import BeamTraversal
from reasoning.answer_extractor import extract

traversal = BeamTraversal(
    adapter=adapter,
    csa_engine=csa,
    embeddings=embeddings,
    communities=community_map,
    beam_width=10,
    max_hop=3,
    # Optional: Bayesian mode (Phase 18)
    probabilistic=True,
    warm_start_strength=5.0,   # Phase 19: reduce cold-start variance
)

paths   = traversal.traverse(["Marie Curie"])
answers = extract(paths, top_k=5)

for ans in answers:
    print(f"{ans.entity_id:30s}  score={ans.score:.4f}")
    print(f"  path: {' -> '.join(ans.best_path.nodes)}")
\end{lstlisting}

---

\subsection{6. LLM Bridge (Optional)}

Convert CEREBRUM reasoning results into prompts or structured context for any LLM:

\begin{lstlisting}
from llm_bridge import generate, to_prompt
from llm_bridge import AnthropicAdapter  # or OpenAIAdapter, OllamaAdapter, HuggingFaceAdapter

# Format as a prompt string
prompt = to_prompt(answers, query="What did Marie Curie discover?")

# Generate an LLM response grounded in the graph paths
result = generate(
    answers=answers,
    query="What did Marie Curie discover?",
    adapter=AnthropicAdapter(api_key=os.environ["ANTHROPIC_API_KEY"], model="claude-sonnet-4-6"),
)

print(result.text)
print(result.grounded_paths)   # The CEREBRUM paths the LLM was given as context
\end{lstlisting}

See \texttt{llm\_bridge/README.md} for adapter confi\-guration details.

---

\subsection{7. Streaming Integration}

\begin{lstlisting}
from adapters.stream_adapter import StreamAdapter
from core.discretizer import STDPDiscretizer, ThresholdDiscretizer
from core.rebalancer import GlobalRebalancer

# Wrap existing adapter with streaming
stream = StreamAdapter(
    base_adapter=adapter,
    window_size=10_000,
    batch_size=50,
)

# Add discretizers
stream.add_discretizer(STDPDiscretizer(
    w_threshold=0.5, n_min=5,
    min_causal_span=1.0,   # adversarial hardening
))

# Background rebalancer
rebalancer = GlobalRebalancer(
    adapter=stream,
    q_drift_threshold=0.05,
    bridge_engine=bridge_engine,   # optional
)
rebalancer.start()

# Push events
stream.push_event("sensor_a", "READS", "temp_42", timestamp=time.time(), weight=0.8)
\end{lstlisting}

---

\subsection{8. Federated Deployment}

\begin{lstlisting}
from adapters.remote_adapter import RemoteAdapter
from adapters.federated_adapter import FederatedAdapter

# Connect to remote CEREBRUM nodes
node_a = RemoteAdapter("https://node-a.example.com", token=os.environ["NODE_A_TOKEN"])
node_b = RemoteAdapter("https://node-b.example.com", token=os.environ["NODE_B_TOKEN"])

# Create federated view
federated = FederatedAdapter([local_adapter, node_a, node_b])

# Query traverses all nodes transparently
traversal = BeamTraversal(adapter=federated, csa_engine=csa, ...)
\end{lstlisting}

---

\subsection{9. Complete Minimal Example}

\begin{lstlisting}
from adapters.csv_adapter import load_csv_adapter
from core.community_engine import best_of_n_dscf
from core.embedding_engine import RandomEngine
from core.attention_engine import CSAEngine
from core.structural_encoder import build_community_distance_matrix, adjacent_community_pairs
from reasoning.traversal import BeamTraversal
from reasoning.answer_extractor import extract

adapter = load_csv_adapter("tests/fixtures/toy_graph.csv")
G       = adapter.to_networkx()

parts   = best_of_n_dscf(G, n_trials=5, seed=42)
cmap    = {node: cid for cid, members in enumerate(parts) for node in members}

emb     = RandomEngine(dim=64).encode_entities({n: n for n in G.nodes()})
csa     = CSAEngine(communities=cmap, embeddings=emb)
csa.set_community_graph(
    build_community_distance_matrix(G, cmap),
    adjacent_community_pairs(G, cmap),
)
csa.set_query_snapshot(cmap)

traversal = BeamTraversal(adapter=adapter, csa_engine=csa, embeddings=emb,
                          communities=cmap, beam_width=10, max_hop=3)
answers   = extract(traversal.traverse(["newton"]), top_k=5)

for a in answers:
    print(f"{a.entity_id:25s}  {a.score:.4f}  {' -> '.join(a.best_path.nodes)}")
\end{lstlisting}

---
\textbf{Copyright © 2026 Bryan Alexander Buchorn. All Rights Reserved.}


\vspace{3em}
\begin{center}
    \small \textit{Project CEREBRUM — March 2026 — Hardened Enterprise Security (v1.2.4)}
\end{center}

\end{document}
